% GENERATED FILE. Edit canonical lessons, then run npm run generate:books.
% canonical-source-hash: fnv1a64:f8022a6d8d7ef58f
% canonical-lessons: PA-W11-form-six-order, PA-W11-form-six-select, PA-W11-form-six-supported, PA-W11-form-six-delayed, PA-W11-form-six-repair, PA-W11-form-six-no-model

\chapter{Join Six Fictional Form Fields}
\label{ch:pa-form-six-field-checkpoint}
\hlchaptermodality{155}

\begin{chapteropening}
\textbf{By the end of this chapter:} \emph{I can select six already taught fictional values from familiar nonverbal cues, write a six-line Gurmukhi form without a visible answer model, and repair selection, spelling, spacing, digit order, and placement separately.}

The learner fills name, language, residence, work, age, and phone lines from remembered cues without romanization or copyable answers, then repairs one error dimension at a time.
\end{chapteropening}

\section[\texorpdfstring{\pa{ਨਾਂ}: · \pa{ਭਾਸ਼ਾ} … (six familiar …)}{ਨਾਂ: · ਭਾਸ਼ਾ … (six familiar …)}]{Six old labels, one gentle order}

\label{lesson:PA-W11-form-six-order}



\begin{quote}
Point to each familiar label. The first four were practised with names and words; the last two with Gurmukhi digits. Date is an already practised seventh field, but it is not part of this six-line rehearsal.
\end{quote}

\subsection*{Writing — labels before values}
Copy only the labels in this order. Leave every value blank. The order is name, language, residence, work, age, phone. Pause after line three.

\begin{quote}
\pa{ਨਾਂ}: \_\_\_\_\_\_
\end{quote}

\begin{quote}
\pa{ਭਾਸ਼ਾ}: \_\_\_\_\_\_
\end{quote}

\begin{quote}
\pa{ਰਿਹਾਇਸ਼}: \_\_\_\_\_\_
\end{quote}

\begin{quote}
\pa{ਕੰਮ}: \_\_\_\_\_\_
\end{quote}

\begin{quote}
\pa{ਉਮਰ}: \_\_\_\_\_\_
\end{quote}

\begin{quote}
\pa{ਫ਼ੋਨ}: \_\_\_\_\_\_
\end{quote}

\subsection*{Your turn — cover and point}
Cover the labels. Point to six empty lines and recall only which label belongs on each. Uncover them and repair one misplaced label, if needed.

\begin{tcolorbox}[breakable,colback=teal!4,colframe=teal!35!black,title={Before you move on}]
Today was an order lesson, not a new vocabulary or exam lesson.
\end{tcolorbox}
\section[\texorpdfstring{A · B · \pa{ਕ} · \pa{ਖ} (old fictional card marks)}{A · B · ਕ · ਖ (old fictional card marks)}]{Reuse the two old card marks}

\label{lesson:PA-W11-form-six-select}



\begin{quote}
The old word cards use A and B. The old digit cards use \pa{ਕ} and \pa{ਖ}. No new name, word, digit, or field is introduced here.
\end{quote}

\subsection*{You'll want to know — point before writing}
Look once at the familiar banks. A chooses \pa{ਅਮਨ}, \pa{ਪੰਜਾਬੀ}, \pa{ਪਿੰਡ}, \pa{ਖੇਤੀ}; \pa{ਕ} chooses \pa{੧੫} and \pa{੦੧੨} \pa{੨੫੧}. B chooses \pa{ਮਨਨ}, \pa{ਹਿੰਦੀ}, \pa{ਸ਼ਹਿਰ}, \pa{ਨੌਕਰੀ}; \pa{ਖ} chooses \pa{੨੫} and \pa{੦੨੫} \pa{੧੨੫}. Point to the matching value in each bank, one line at a time. The paired marks A · \pa{ਕ} or B · \pa{ਖ} choose one fictional person's six old values; they are not a new sound or spelling rule.

\subsection*{Your turn — one card, six choices}
Choose B · \pa{ਖ}. Without writing a form, say which bank side belongs on the first four lines and which side on the final two. Check the mapping above.

\begin{tcolorbox}[breakable,colback=teal!4,colframe=teal!35!black,title={Before you move on}]
The next lesson keeps these banks visible while the six lines are written.
\end{tcolorbox}
\section[\texorpdfstring{\pa{ਨਾਂ}: \pa{ਅਮਨ} · \pa{ਭਾਸ਼ਾ} … (supported …)}{ਨਾਂ: ਅਮਨ · ਭਾਸ਼ਾ … (supported …)}]{Six lines with the old banks visible}

\label{lesson:PA-W11-form-six-supported}



\begin{quote}
Return to the six blank labels. Point to the first three lines, pause, then point to the last three. Today the old A and \pa{ਕ} values stay visible.
\end{quote}

\subsection*{Writing — supported entry}
For A · \pa{ਕ}, use the old visible word bank: \textbf{\pa{ਅਮਨ} · \pa{ਪੰਜਾਬੀ} · \pa{ਪਿੰਡ} · \pa{ਖੇਤੀ}}. Use the old visible digit bank: \textbf{\pa{੧੫} · \pa{੦੧੨} \pa{੨੫੧}}. Copy one selected value beside each familiar label in the fixed order. Stop after the first three lines and check only which value belongs to which label; then do the rest.

\begin{quote}
\pa{ਨਾਂ}: \_\_\_\_\_\_
\end{quote}

\begin{quote}
\pa{ਭਾਸ਼ਾ}: \_\_\_\_\_\_
\end{quote}

\begin{quote}
\pa{ਰਿਹਾਇਸ਼}: \_\_\_\_\_\_
\end{quote}

\begin{quote}
\pa{ਕੰਮ}: \_\_\_\_\_\_
\end{quote}

\begin{quote}
\pa{ਉਮਰ}: \_\_\_\_\_\_
\end{quote}

\begin{quote}
\pa{ਫ਼ੋਨ}: \_\_\_\_\_\_
\end{quote}

\subsection*{Your turn — one boundary at a time}
Check each colon has one readable space after it. The phone number also keeps its old single group space. This is visible-bank practice, not independent writing evidence.

\begin{tcolorbox}[breakable,colback=teal!4,colframe=teal!35!black,title={Before you move on}]
Next time the banks will be covered before writing.
\end{tcolorbox}
\section[\texorpdfstring{\pa{ਨਾਂ}: · \pa{ਭਾਸ਼ਾ} … (delayed …)}{ਨਾਂ: · ਭਾਸ਼ਾ … (delayed …)}]{Cover the six old values}

\label{lesson:PA-W11-form-six-delayed}



\begin{quote}
Look once at the old B and \pa{ਖ} values: \textbf{\pa{ਮਨਨ} · \pa{ਹਿੰਦੀ} · \pa{ਸ਼ਹਿਰ} · \pa{ਨੌਕਰੀ}}, then \textbf{\pa{੨੫} · \pa{੦੨੫} \pa{੧੨੫}}. Cover this complete bank and wait ten seconds.
\end{quote}

\subsection*{Writing — delayed entry}
With the bank still hidden, fill the six lines for B · \pa{ਖ}. Recall one value at a time. Pause after the residence line. Do not uncover the bank until all six attempts are complete.

\begin{quote}
\pa{ਨਾਂ}: \_\_\_\_\_\_
\end{quote}

\begin{quote}
\pa{ਭਾਸ਼ਾ}: \_\_\_\_\_\_
\end{quote}

\begin{quote}
\pa{ਰਿਹਾਇਸ਼}: \_\_\_\_\_\_
\end{quote}

\begin{quote}
\pa{ਕੰਮ}: \_\_\_\_\_\_
\end{quote}

\begin{quote}
\pa{ਉਮਰ}: \_\_\_\_\_\_
\end{quote}

\begin{quote}
\pa{ਫ਼ੋਨ}: \_\_\_\_\_\_
\end{quote}

\subsection*{Your turn — compare after the attempt}
Uncover the old bank. Compare only the first differing line. Check its selection, then spelling or digit order, then the space after the colon. Correct that line without erasing lines that were already right.

\begin{tcolorbox}[breakable,colback=teal!4,colframe=teal!35!black,title={Before you move on}]
This is delayed practice. The later checkpoint will remove the bank entirely.
\end{tcolorbox}
\section[\texorpdfstring{\pa{ਰਿਹਾਇਸ਼}: \pa{ਸ਼ਹਿਰ} … (repair six …)}{ਰਿਹਾਇਸ਼: ਸ਼ਹਿਰ … (repair six …)}]{Repair one dimension at a time}

\label{lesson:PA-W11-form-six-repair}



\begin{quote}
Read the six old labels in their new order. Today we repair old forms, not learn a seventh field or a new sentence pattern.
\end{quote}

\subsection*{Writing — six small repair passes}
Use B · \pa{ਖ}. First check that each line has the selected B or \pa{ਖ} value. Next check spelling in the four word lines; \textbf{\pa{ਸ਼ਹਰ}} needs the old vowel mark to become \textbf{\pa{ਸ਼ਹਿਰ}}. Next check one space after every colon. Then check the two Gurmukhi digit lines left to right and the phone's one group space. Last, check that each whole value sits beside its own label. Change only the first wrong dimension you find on each pass.

The work values here are the already taught nouns \textbf{\pa{ਖੇਤੀ}} and \textbf{\pa{ਨੌਕਰੀ}}; there is no agreement-bearing sentence to score in this form. Agreement remains a separate task when a sentence actually requires it.

\subsection*{Your turn — one visible error}
For B · \pa{ਖ}, repair only the last two digit positions in \textbf{\pa{ਫ਼ੋਨ}: \pa{੦੨੫} \pa{੧੫੨}}. Keep the label, colon, space, first digit group, and other form lines as they were.

\begin{tcolorbox}[breakable,colback=teal!4,colframe=teal!35!black,title={Before you move on}]
The next lesson asks for an untimed six-line attempt without the bank.
\end{tcolorbox}
\section[\texorpdfstring{\pa{ਨਾਂ}: · \pa{ਭਾਸ਼ਾ} … (independent …)}{ਨਾਂ: · ਭਾਸ਼ਾ … (independent …)}]{One untimed six-field form}

\label{lesson:PA-W11-form-six-no-model}



\begin{quote}
Close the earlier lessons. No value bank, support-language value label, romanization, Latin-digit version, or copyable Gurmukhi answer appears below. Use only the old fictional card marks B · \pa{ਖ}.
\end{quote}

\subsection*{Writing — independent controlled choice}
Write all six requested values from memory. Pause after the third line. Keep the last two values in Gurmukhi digits. Do not reopen an earlier lesson or use a dictionary, translator, or spell-checker until the attempt is complete.

\begin{quote}
B · \pa{ਖ}
\end{quote}

\begin{quote}
\pa{ਨਾਂ}: \_\_\_\_\_\_
\end{quote}

\begin{quote}
\pa{ਭਾਸ਼ਾ}: \_\_\_\_\_\_
\end{quote}

\begin{quote}
\pa{ਰਿਹਾਇਸ਼}: \_\_\_\_\_\_
\end{quote}

\begin{quote}
\pa{ਕੰਮ}: \_\_\_\_\_\_
\end{quote}

\begin{quote}
\pa{ਉਮਰ}: \_\_\_\_\_\_
\end{quote}

\begin{quote}
\pa{ਫ਼ੋਨ}: \_\_\_\_\_\_
\end{quote}

\subsection*{Your turn — repair after the attempt}
Only after all six lines are filled, reopen the old banks. Check selection, word spelling, colon spacing, digit order and phone grouping, then field placement in separate passes. Repair one differing dimension at a time.

\begin{tcolorbox}[breakable,colback=teal!4,colframe=teal!35!black,title={Before you move on}]
This is an untimed minimum-size form rehearsal, not a timed A1 paper. Date remains a separately practised seventh field. A named-reader message, timed production, full mocks, and human A1 validation remain separate work.
\end{tcolorbox}
